\documentclass{article}
\input{~/studies/preamble}

\title{Autoencoder}
\author{Hyoung Joon, Kim}
\date{\today}

\begin{document}
\maketitle
\tableofcontents
\newpage

\section{Introduction}
    \textbf{Autoencoder} produces outputs that are similar to input data. It first maps the input $x$ to
    a hidden representation $z(x)$ and produce the output $y(z)$ which reserves the original data $x$ well.
    This can be seen as two parts, \emph{encoder} and \emph{decoder}. Encoder can later be used at dimensionality reduction.
    To inhibit autoendcoder from generating trivial solution, we give some constraints, such as restricting the
    dimension of hidden representation or make it to have sparse representation.

\section{Deterministic Autoencoder}
    \subsection{Deep Autoencoder}
        \begin{defn}{Deep Autoencoder}{deepAutoenc}
            Let $\{\tb x_n\} \in \mathbb{R}^D$ be a given data, and let $\tb{y}(\tb x, \tb w) : \mathbb{R}^D \to \mathbb{R}^D$
            be a deep network function with parameter $\tb w$ having bottleneck.
            Then, the \textbf{deep autoencoder} minizes the reconstruction error:
            \[
            E(\tb w)=\sum_{n=1}^{N}\|\tb{y}(\tb x_n, \tb w) - \tb x_n\|^2
            \]
        \end{defn}
        Usually the network $\tb y(\tb x, \tb w)$ decrease the dimensionality and then increase dimensionality of hidden units to stop autoencoder from learning
        trivial representations, which is called bottleneck.
    \subsection{Sparse Autoencoder}
        Instead of making bottleneck like deep autoencoder, we can use regularizer so that the internal representation have a sparse representation.
        This method have an effect of inducing the effective dimensionality of internal representation.
        One simple choice of sparse regularizer is L1 regularizer, 
        \begin{defn}{Sparse Autoencoder with L1 Regularizer}{sparseAutoenc}
            Let $\{\tb x_n\} \in \mathbb{R}^D$ be a given data, and let $\tb{y}(\tb x, \tb w) : \mathbb{R}^D \to \mathbb{R}^D$
            be a deep network function with parameter $\tb w$.
            Then, the \textbf{sparse autoencoder} minimizes the reconstruction error with L1 regularizer:
            \[
            E(\tb w)=\sum_{n=1}^{N}\|\tb{y}(\tb x_n, \tb w) - \tb x_n\|^2 + \lambda\sum_{k=1}^{K}|z_k|
            \]
            where $z_k$ are the appropriate hidden units according to the internal representation.
        \end{defn}
    \subsection{Denoising Autoencoder}
        Another constraint is to corrupt the given data. Then the autoencoder will have ability to denoise the corrupted data.
        \begin{defn}{Denoising Autoencoder}{denoisAutoenc}
            Let $\{\tb x_n\} \in \mathbb{R}^D$ be a given data, and let $\tb{y}(\tb x, \tb w) : \mathbb{R}^D \to \mathbb{R}^D$
            be a deep network function with parameter $\tb w$.
            Then, the \textbf{denoising autoencoder} minizes the following error function:
            \[
            E(\tb w)=\sum_{n=1}^{N}\|\tb{y}(\tilde{\tb x}_n, \tb w) - \tb x_n\|^2
            \]
            where $\tilde{\tb x}_n$ are corrupted dataset. 
        \end{defn}
        The denoising autoencoder's training is related to score mathcing where the score is defined as $\tb s = \nabla_{\tb x}\log{p(\tb x)}$.
        The autoencoder learns to reverse the corruption, and the score vector $\nabla\log{p(x)}$ points towards high data density,
        These two methods are similar and have deep relationship.
    \subsection{Masked Autoencoder}
        The transformer models such as BERT are trained by masking the random subsets of input texts, and this can be applied to autoencoder models, too.
        This is usually combined with transformer models to create a vision transformer architecture.
        Then the autoencoder trains to predict the deleted pixels, inferencing with leftover pixels.

\section{Variational Autoencoder}
    \subsection{Introduction}
        The \textbf{variational autoencoder}, unlike deterministic autoencoder, does not fix the internal representation, but allow some uncertainty
        of internal representation. However, the likelihood defined by a neural network make the training and inference intractable.
        Variational autoencoder resolve this problem by using \emph{variational inference} technique, approximating the likelihood function with a
        feasible function.
    \subsection{Amortized Inference}
        Consider a generative model with latent variable having distribution $p(\tb z)=\mathcal{N}(\tb z|\tb 0, \tb I)$.
        We can make a network that creates data from latent variables: $\tb g(\tb z,\tb w)$. And we can model a probability distribution with network $\tb g$,
        for example, a Gaussian distribution $p(\tb x|\tb z, \tb w) = \mathcal{N}(\tb x|\tb g(\tb z,\tb w), \tb I)$.
        
        Now we derive the ELBO of variational autoencoder.
        The log-likelihood of data $\tb x$ is:
        \[
        \log{p(\tb x| \tb w)=\mathcal{L}(\tb w)+\text{KL}(q(\tb z )\|p(\tb z|\tb x, \tb w))}
        \]
        where $\mathcal{L}(\tb w)$ is the evidence lower bound,
        \[
        \mathcal{L}(\tb w)=\int q(\tb z)\log{\left\{\frac{p(\tb x|\tb z, \tb w)p(\tb z)}{q(\tb z)}\right\}}d\tb z
        \]
        and
        \[
        \text{KL}(q(\tb z) \| p(\tb z|\tb x,\tb w))=-\int q(\tb z)\log{\left\{ \frac{p(\tb z|\tb x,\tb w)}{q(\tb z)} \right\}}d\tb z.
        \]
        The difference from GMM or pPCA is that the E step is infeasible, since the exact posterior distribution requires the
        likelihood $p(\tb x|\tb w)$, which is intractable:
        \[
        p(\tb z|\tb x, \tb w)=\frac{p(\tb x|\tb z, \tb w)p(\tb z)}{p(\tb x|\tb w)}
        \]
        The solution is to approximate the posterior distribution $p(\tb z|\tb x, \tb w)$,
        with another network $q(\tb z|\tb x, \bm\phi)$, called the \emph{encoder network}.
        The original network then can be called as \emph{decoder network}.

        A typical choice for the encoder network is a Guassian ditribution with mean and variance given by the function that takes $\tb x$ as input;
        \[
        q(\tb z|\tb x,\bm\phi)=\prod_{j=1}^{M}\mathcal{N}(z_j|\mu_j(\tb x,\bm\phi), \sigma^2_j(\tb x,\bm\phi))
        \]
        The E step now coressponds to updating $\bm\phi$ and M step corresponds to updating $\tb w$. Notice that the E step does not in general reduce the
        KL divergence to zero, since it is just an approximation.
    \subsection{Reparametrization Trick}
        Now, we introduce the \textbf{reparametrization trick}. The lower bound $\mathcal{L}(\tb w)$ is still intractable since it involves integrals over the latent
        variables $\{\tb z_n\}$:
        \begin{align*}
            \mathcal{L}(\tb w,\bm\phi)&=\int q(\tb z|\tb x, \bm\phi)\log{\left\{ \frac{p(\tb x|\tb z, \tb w)p(\tb z)}{q(\tb z|\tb x, \bm\phi)} \right\}}d\tb z
            \\ &=\int q(\tb z|\tb x,\bm\phi)\log{p(\tb x|\tb z,\tb w)}d\tb z - \text{KL}(q(\tb z|\tb x,\bm\phi)\|p(\tb z))
        \end{align*}
        The first term is called \emph{reconstruction error}, since it can be seen as expectation of likelihood of $x$ about latent variable created from given data $x$.
        The second term, is called regularization term, and can be evaluated analytically:
        \[
        \text{KL}(q(\tb z|\tb x,\bm\phi)\|p(\tb z))=\frac{1}{2}\sum_{j=1}^{M}\{1+\log{\sigma^2_j(\tb x, \bm\phi)}-\mu^2_j(\tb x,\bm\phi)-\sigma^2_j(\tb x,\bm\phi)\}
        \]
        The leftover term can be approximated by simple Monte Carlo estimator:
        \[
        \int q(\tb z|\tb x,\bm\phi)\log{p(\tb x|\tb z,\tb w)}d\tb z \simeq \frac{1}{L}\sum_{l=1}^{L}\log{p(\tb x|\tb z^{(l)}, \tb w)}
        \]
        where $\{\tb z^{(l)}\}$ are samples drawn from the encoder distribution $q(\tb z|\tb x, \bm\phi)$. The problem is that the samples are fixed.
        The change in parameter $\bm\phi$ will change the distribution where the samples are sampled from, but the sample stays fixed.
        To resolve this, we use reparametrization trick. Instead of sampling directly from $q(\tb z|\tb x,\bm\phi)$, we use
        that $\tb z_j = \sigma_j\epsilon + \mu_j$ where $\epsilon \sim \mathcal{N}(0,1)$ has Guassian distribution with
        mean $\mu_j$ and variance $\sigma_j^2$, and sample $\epsilon \sim \mathcal{N}(0,1)$ to create samples $\tb z^{(l)}$.
        This makes the dependence on $\bm\phi$ explicit and allow us to use gradient descent methods.
    \subsection{Algorithm}
        The training of VAE must be balanced over two terms, the reconstruction error and KL divergence.
        
        \begin{algorithm}[H]\label{alg:VAEtrain}
            \DontPrintSemicolon
            \SetKwInOut{Input}{Input}
            \SetKwInOut{Output}{Output}
            \SetArgSty{textnormal}
            
            \caption{Variational autoencoder training algorithm}

            \Input{dataset $\data = \{\tb x_1,\dots,\tb x_N\}$ \newline
            encoder network $q(\tb z|\tb x,\bm\phi)=\prod_{j=1}^{M}\mathcal{N}(z_j|\mu_j(\tb x,\bm\phi),\sigma^2_j(\tb x,\bm\phi))$ \newline
            decoder network $\tb g(\tb z, \tb w)$ \newline
            initial weight vectors $\tb w$, $\bm\phi$ \newline
            learning rate $\eta$}
            \Output{final weight vectors $\tb w$, $\bm\phi$}

            \BlankLine
            \hrule

            \Repeat{convergence}
            {
                \For{$n \in$ mini-batches of data}
                {
                    $\mathcal{L} \leftarrow 0$\;
                    \For{$j \in \{1,\dots,M\}$}
                    {
                        $\epsilon_{nj} \leftarrow \mathcal{N}(0,1)$\;
                        $z_{nj} \leftarrow \mu_j(\tb x_n, \bm\phi) + \sigma_j(\tb x_n,\bm\phi)\epsilon_{nj}$\;
                        $\mathcal{L} \leftarrow \mathcal{L} + \frac{1}{2}\{1+\log{\sigma^2_j(\tb x_n,\bm\phi)-\mu_j^2(\tb x_n,\bm\phi)-\sigma^2_j(\tb x_n,\bm\phi)}\}$\;
                    }
                    $\mathcal{L} \leftarrow \mathcal{L} + \log{p(\tb x_n|\tb z_n,\tb w)}$\;
                    $\tb w \leftarrow \tb w + \eta\nabla_{\tb w}\mathcal{L}$\;
                    $\bm\phi \leftarrow \bm\phi + \eta\nabla_{\bm\phi}\mathcal{L}$\;
                }
            }
            \Return{$\tb w$, $\bm\phi$}
        \end{algorithm}
        The decoder network $\tb g(\tb z, \tb w)$ is inherent in $p(\tb x|\tb z,\tb w)$. For example, when the variables are continuous, $p(\tb x|\tb z, \tb w)=\mathcal{N}(\tb g(\tb z, \tb w), \tb I)$
        
        \subsubsection{Posterior Collapse and Holes in Latent Space}
            The VAE training algorithm, Algorithm \ref{alg:VAEtrain} has a problem called \textbf{posterior collapse}.
            When the variational distribution $q(\tb z|\tb x, \tb w)$ converges to the prior distribution $p(z)$ and ignores $\tb x$, leading to poor reconstruction result.
            This happens because the second term in equation
            \begin{align*}
                \mathcal{L}(\tb w,\bm\phi)&=\int q(\tb z|\tb x, \bm\phi)\log{\left\{ \frac{p(\tb x|\tb z, \tb w)p(\tb z)}{q(\tb z|\tb x, \bm\phi)} \right\}}d\tb z
                \\ &=\int q(\tb z|\tb x,\bm\phi)\log{p(\tb x|\tb z,\tb w)}d\tb z - \text{KL}(q(\tb z|\tb x,\bm\phi)\|p(\tb z)),
            \end{align*}
            the KL divergence term, encourage $q(\tb z|\tb x, \tb w)$ to be close to the prior $p(\tb z)$. Thus the posterior collapse can happen when the KL divergence term is relatively closer to zero 
            than the reconstruction term.
            
            Different problem can happen when the latent code is not compressed well. When the latent code is not compressed and therefore is very different from prior,
            the reconstruction result may be highly accurate, but the samples from prior won't generate realistic outupts. We say that there are \textbf{holes in latent space}.
            In this case, the KL divergence term would be relatively larger than the reconstruction error.
        
        \subsubsection{Solution}
            We can address both problems by multiplying a coeffiecient $\beta$ in front of reconstruction term, to control the regularization effectiveness of KL divergence term.
            Typically, $\beta > 1$. If the reconstruction is poor, we increse $\beta$. When the samples are poor, we decrease $\beta$. 
            Or we can change $\beta$ during training session, by starting with a small value and gradually increasing.

            Another method is to modify the KL divergence term. Wasserstein Autoencoder, VQ-VAE and InfoVAE is the examples of it.
            
            \emph{Guess I should study these methods when I build VAE from scratch...}
\end{document}